% !TEX root = ../jlt.tex
\section{Related Work}
\label{sec:related}

\paragraph{Type classes.}
Dictionary passing is Wadler and Blott's~\cite{wadler1989adhoc}.
Table~\ref{tab:dictionary} is that translation with \Cpp{}'s names in the left column, and the row without a counterpart there, structural satisfaction, is a difference Bernardy et al.~\cite{bernardy2010generic} weigh in their comparison of concepts and type classes.
Coherence, that a program has one meaning however its constraints are resolved, is where the designs diverge: Racordon et al.~\cite{racordon2025coherence} survey how Haskell, Swift, Rust and Scala keep or trade it, between unique instances and context-dependent resolution.
FJlt takes the strict side by construction, since an instance is a function specialisation and the one-definition rule admits one per head.
Theorem~\ref{thm:coherence} is that choice stated as a property of the translation, and the price, that instances for families of types declared after their shape do not fit, is in Section~\ref{sec:limits}.
Instance resolution is the inference half of type classes, Hindley--Milner's~\cite{milner1978theory,damas1982principal} extended with predicates~\cite{jones1994qualified}; FJlt infers types only locally~\cite{pierce2000local}, a result type from its body, but it infers laws the same way.
Rust and Swift also check a generic body against its bounds: a comparison on an unconstrained type parameter is a compile error in both, and Rust's message names the missing bound~\cite{klabnik2023rust,swift2026generics}.
That is Rule~2 as a compiler check, the half \Cpp{20} dropped, and the posture, together with the linter of Section~\ref{sec:limits}, is what recovering it by hand costs.

\paragraph{Rewrite rules and fusion.}
GHC lets a library declare equations the optimiser rewrites by~\cite{peytonjones2001rules}, wherever their pattern matches and unchecked.
A law in Jlt is as unchecked, a certificate the compiler finds but does not prove (Section~\ref{sec:limits}); the difference is where it lives and what it does: a declared instance, so what a program assumes is in its types, licensing one overload that resolution chooses once.
Fusion removes the lists a pipeline would build between its stages: Wadler's deforestation by unfolding the definitions~\cite{wadler1990deforestation}, shortcut fusion by one rule on \texttt{foldr} and \texttt{build}~\cite{gill1993shortcut}, stream fusion by rewriting the stages to a stream and inlining~\cite{coutts2007stream}.
A pipeline of \Cpp{} ranges builds none: a range already is the stream, so inlining alone fuses Listing~\ref{lst:sites}'s stages into one loop (Table~\ref{tab:binding}), and the posture adds nothing there.

\paragraph{Concepts, and the calculus next door.}
The nearest neighbour of FJlt is Siek and Lumsdaine's $F^G$~\cite{siek2005essential,siek2004formalization}: System~F with concepts and models, translated into System~F by dictionary passing, with the translation proved sound in Isabelle.
It grew from the same work as the concepts then proposed for \Cpp{0x}~\cite{gregor2006concepts,dosreis2006specifying}, which checked template bodies and were removed~\cite{stroustrup2009remove}, leaving the gap Rule~2 names.
$F^G$'s models are nominal and can be declared for whole families of types, which is what the library behind this paper's listings needs after its shapes, where FJlt's function-specialisation instances are not found (Section~\ref{sec:limits}).
FJlt is the smaller calculus, and the price of being a subset of a shipped language.
Of the generic-programming features Garcia et al.~\cite{garcia2007extended} compare across eight languages, two that \Cpp{} then had no explicit support for, multi-type concepts and retroactive modelling, are what Listing~\ref{lst:rule2} exercises.

\paragraph{Featherweight calculi.}
Featherweight Java~\cite{igarashi2001featherweight} set the pattern: a core calculus small enough to prove things about, with a syntax that is a subset of the host language's.
Featherweight Go~\cite{griesemer2020featherweight} follows the pattern for structural subtyping and monomorphisation, and \emph{Generic Go to Go}~\cite{ellis2022generic} proves a dictionary-passing translation correct by a bisimulation.
That is the genre Section~\ref{sec:guarantees} is in, preservation and simulation of a translation, with soundness, termination and parametricity inherited from the target rather than proved from scratch.

\paragraph{Formal semantics of \Cpp{}.}
CoreC++~\cite{wasserrab2006operational} gives an operational semantics and a machine-checked type safety proof for multiple inheritance; Ramananandro, Dos Reis and Leroy~\cite{ramananandro2011layout,ramananandro2012construction} mechanise object layout, construction and destruction in Coq.
These formalise the part of \Cpp{} that Rule~1 excludes, objects with identity and lifetime, and none covers templates, concepts or constant evaluation, so FJlt is a fragment with its own grammar rather than a subset of a formalised \Cpp{}.

\paragraph{What System~F supplies.}
Strong normalisation is Girard's \cite{girard1989proofs}, with inductive types Matthes' \cite{matthes1999extensions}, and parametricity Reynolds' \cite{reynolds1983types}: the appendix uses both as black boxes, which is the point of translating into System~F; data types as their own folds is B\"ohm and Berarducci's reading \cite{bohm1985automatic}, which the translation takes with iterators in place of Church encodings.

\paragraph{Functional programming in \Cpp{}, and the other direction.}
Functional programming in \Cpp{} is practised and taught: in a book~\cite{cukic2018functional}, in a metaprogramming library~\cite{chagnon2016hana}, and in category theory with examples in \Cpp{}~\cite{milewski2019category}; none states a fragment or proves anything of it.
The opposite route to the same end verifies a program in a proof assistant and extracts C from it, as Low* does from F*~\cite{protzenko2017verified}; Jlt keeps the programmer in \Cpp{} and puts the proof in the translation.

\paragraph{Coding standards and profiles.}
The documents an audit of \Cpp{} works from are rule catalogues: MISRA \Cpp{}~\cite{misra2023cpp} and the SEI CERT standard~\cite{ballman2016cert} list hundreds of rules between them, each justified by the defects it prevents, with no statement of what a compliant program computes.
Jlt is narrower, a pure computational core with no I/O, and so complements them rather than replacing them; where they meet they can differ, as MISRA forbids recursion except in constant evaluation and Rule~3 admits structural recursion on the strength of a termination argument.
The nearest relatives are the safety profiles, now aimed at \Cpp{29}~\cite{sutter2025profiles,sutter2026london}, named sets of rules the compiler would enforce, and the criticism that a profile should come with the invariant it upholds~\cite{muller2025principled}: Jlt reads as a functional profile, held by review and a linter rather than wholly by the compiler, and its three rules state theirs.

